\section{Cotangent, five-part, Cagnoli, right, and quadrantal rules}
\label{sec:cotangent-special}

\subsection{The six cotangent four-part formulas}

The six parts of a spherical triangle occur cyclically as
\(a,C,b,A,c,B\).  Choosing four consecutive parts gives one inner side, one
inner angle, one outer side and one outer angle.  The general mnemonic is
\begin{equation}
 \cos(\text{inner side})\cos(\text{inner angle})
 =\cot(\text{outer side})\sin(\text{inner side})
  -\cot(\text{outer angle})\sin(\text{inner angle}).        \tag{8.1}
\end{equation}
Explicitly,
\begin{align}
 \cos b\cos C&=\cot a\sin b-\cot A\sin C,                 \tag{CT1}\\
 \cos b\cos A&=\cot c\sin b-\cot C\sin A,                 \tag{CT2}\\
 \cos c\cos A&=\cot b\sin c-\cot B\sin A,                 \tag{CT3}\\
 \cos c\cos B&=\cot a\sin c-\cot A\sin B,                 \tag{CT4}\\
 \cos a\cos B&=\cot c\sin a-\cot C\sin B,                 \tag{CT5}\\
 \cos a\cos C&=\cot b\sin a-\cot B\sin C.                 \tag{CT6}
\end{align}
For a proper triangle every side and angle sine is positive, so all cotangents
have canonical apartness witnesses.

\begin{proof}[Proof of CT1 and the family]
Begin with the cosine law for \(a\), substitute the cosine law for \(c\), and
use the sine law to replace \(\sin c\):
\begin{align*}
 \cos a
 &=\cos b\cos c+\sin b\sin c\cos A\\
 &=\cos b(\cos a\cos b+\sin a\sin b\cos C)
   +\sin b\sin C\sin a\cot A.
\end{align*}
Move the \(\cos a\cos^2b\) term to the left:
\[
 \cos a\sin^2b
 =\cos b\sin a\sin b\cos C
  +\sin b\sin C\sin a\cot A.
\]
Cancel the positive product \(\sin a\sin b\) to obtain CT1.  The same
substitution pattern gives CT3 and CT5.  Applying those three identities to
the polar triangle gives CT2, CT4 and CT6.  Before cancellation, each proof is
an integer polynomial identity, so it is valid without any division.
\end{proof}

Two frequently printed consequences are already visible here:
\begin{align}
 \cos A&=\csc b\,\csc c\,(\cos a-\cos b\cos c),           \tag{8.2}\\
 \cos a\cos C&=\sin a\cot b-\sin C\cot B.                 \tag{8.3}
\end{align}

\subsection{Five-part rules}

Substitute the cosine law for \(b\) into the law for \(a\):
\begin{align*}
 \cos a
 &=(\cos a\cos c+\sin a\sin c\cos B)\cos c
   +\sin b\sin c\cos A.
\end{align*}
After moving \(\cos a\cos^2c\) to the left and canceling the positive factor
\(\sin c\), one obtains
\begin{equation}
 \cos a\sin c
 =\sin a\cos c\cos B+\sin b\cos A.                         \tag{8.4}
\end{equation}
Equivalently,
\begin{equation}
 \sin a\cos B=\cos b\sin c-\sin b\cos c\cos A.            \tag{8.5}
\end{equation}
Cyclic permutation gives
\begin{align}
 \cos b\sin a&=\sin b\cos a\cos C+\sin c\cos B,          \tag{8.6}\\
 \cos c\sin b&=\sin c\cos b\cos A+\sin a\cos C.          \tag{8.7}
\end{align}
Polar duality and reversal of the cyclic order produce the remaining standard
five-part variants.  Their division-free versions require no side conditions.

\subsection{Cagnoli's equation}

\begin{theorem}[Cagnoli]
\label{thm:cagnoli}
The six parts of a spherical triangle satisfy
\begin{equation}
 \sin b\sin c+\cos b\cos c\cos A
 =\sin B\sin C-\cos B\cos C\cos a.                         \tag{8.8}
\end{equation}
There are two cyclic companions.
\end{theorem}

\begin{proof}
Multiply the side cosine law for \(a\) by \(\cos A\), replacing
\(\cos^2A\) by \(1-\sin^2A\):
\begin{equation}
 \cos a\cos A
 =\cos b\cos c\cos A+\sin b\sin c
  -\sin b\sin c\sin^2A.                                   \tag{8.9}
\end{equation}
Multiply the angle cosine law for \(A\) by \(\cos a\):
\begin{equation}
 \cos a\cos A
 =-\cos B\cos C\cos a+\sin B\sin C
  -\sin B\sin C\sin^2a.                                  \tag{8.10}
\end{equation}
The sine law gives
\(
 \sin b\sin c\sin^2A=\sin B\sin C\sin^2a
\).
Subtract the common remainder from (8.9)--(8.10), yielding (8.8).
\end{proof}

\subsection{Right spherical triangles}

Assume \(C=\pi/2\).  Then \(\sin C=1\) and \(\cos C=0\).  The core laws reduce
to Napier's ten rules:
\begin{alignat}{3}
 \text{(R1)}\quad&\cos c=\cos a\cos b,
 &\qquad \text{(R6)}\quad&\tan b=\cos A\tan c,             \tag{8.11}\\
 \text{(R2)}\quad&\sin a=\sin A\sin c,
 &\text{(R7)}\quad&\tan a=\cos B\tan c,                  \tag{8.12}\\
 \text{(R3)}\quad&\sin b=\sin B\sin c,
 &\text{(R8)}\quad&\cos A=\sin B\cos a,                  \tag{8.13}\\
 \text{(R4)}\quad&\tan a=\tan A\sin b,
 &\text{(R9)}\quad&\cos B=\sin A\cos b,                  \tag{8.14}\\
 \text{(R5)}\quad&\tan b=\tan B\sin a,
 &\text{(R10)}\quad&\cos c=\cot A\cot B.                 \tag{8.15}
\end{alignat}
The tangent and cotangent equations are guarded by the corresponding cosine or
sine apartness witnesses.  R1 is the side cosine law; R2--R3 are the sine law;
R8--R9 come from the angle cosine laws.  R4--R7 and R10 follow by guarded
division.  Thus the ten-rule mnemonic introduces no new theorem.

Napier's circle summarizes them constructively.  Write the five parts
\[
 a,\quad \frac\pi2-B,\quad b,\quad \frac\pi2-c,\quad \frac\pi2-A
\]
in cyclic order.  For any three selected parts, the sine of the middle part is
the product of the tangents of the two adjacent parts and also the product of
the cosines of the two opposite parts.  Each assertion is one of
(R1)--(R10) after the complement identities.

\subsection{Quadrantal spherical triangles}

Assume instead that \(c=\pi/2\).  Applying the right-triangle rules to the
polar triangle gives
\begin{alignat}{3}
 \text{(Q1)}\quad&\cos C=-\cos A\cos B,
 &\qquad \text{(Q6)}\quad&\tan B=-\cos a\tan C,            \tag{8.16}\\
 \text{(Q2)}\quad&\sin A=\sin a\sin C,
 &\text{(Q7)}\quad&\tan A=-\cos b\tan C,                   \tag{8.17}\\
 \text{(Q3)}\quad&\sin B=\sin b\sin C,
 &\text{(Q8)}\quad&\cos a=\sin b\cos A,                  \tag{8.18}\\
 \text{(Q4)}\quad&\tan A=\tan a\sin B,
 &\text{(Q9)}\quad&\cos b=\sin a\cos B,                  \tag{8.19}\\
 \text{(Q5)}\quad&\tan B=\tan b\sin A,
 &\text{(Q10)}\quad&\cos C=-\cot a\cot b.                \tag{8.20}
\end{alignat}
Again the quotient forms are derived only under apartness.  The minus signs in
Q1, Q6, Q7 and Q10 arise from
\(\cos(\pi-x)=-\cos x\) under polar duality.
